% Methodology section (3.1-3.9) of the journal paper; \input by main.tex.
%
% Content from the final report (Final_correct_ results/IMD_Mapping_report 1.pdf,
% 2026-10-23, Sections 2.4 and 3), checked against gee_models.py (Earth Engine
% smileRandomForest in REGRESSION mode) and 05_Independent_Validation.ipynb
% (levels = value rounded to the nearest ninth; binary threshold >= 50%).
% Figures are \figfrom placeholders naming the final-report figure to copy in.

\section{Methodology}
\label{sec:method}

\begin{figure}[htbp]
  \centering
  \figfrom{final report Fig.~6 --- workflow, five steps and one model}
  \caption{Workflow. Training targets (CLMS, GHSL) and predictors (AlphaEarth
  embeddings, Sentinel-2 composites) enter at sampling and extraction; a
  random forest is tuned, trained and applied in Google Earth Engine; the
  EarthLabel plots enter only at the final, independent validation.}
  \label{fig:workflow}
\end{figure}

The workflow (\cref{fig:workflow}) uses one model, a random forest
\cite{breiman2001rf}, and the same sequence of steps for every city,
training target and predictor set: stratified sampling inside a common valid
mask, extraction of the predictors at the sample points, tuning by spatially
blocked cross-validation, prediction of the full 10\,m map, and validation
against independent photo-interpreted plots. Tuning, training and
prediction all run in Google Earth Engine \cite{gorelick2017gee} through its
Python API, so each map is produced by exactly the model that was tuned. The
code is openly available, and the predicted rasters and input data are
deposited on Zenodo \cite{brovelli2026zenodo}.

% ---------------------------------------------------------------------------
\subsection{Training-sample generation}
\label{sec:sampling}

Sample points must cover the study area, represent the full range of
imperviousness, and fall where every predictor has a value.

\paragraph{Valid mask.} A pixel is valid when it is not water in ESA
WorldCover \cite{zanaga2022worldcover} and every predictor of the city has a
value there: the AlphaEarth embedding, the Sentinel-2 median (and, in Milan,
percentile) composite, and the Sentinel-2 stack. The stack keeps each date
as separate bands and is not gap-filled (\cref{sec:compositing}), so a pixel
that is cloud or shadow on any stack date has no stack value. All points
are drawn from valid pixels only, so every predictor has a value at every
point and all models of a city are trained on the same points.
\Cref{tab:mask} gives the share of each study area removed by the mask:
water accounts for most of it, the stack for the remainder in Milan and
Hanoi, and the median composite and AlphaEarth remove nothing.

\begin{table}[htbp]
  \centering
  \small
  \caption{Share of each study area removed by the valid mask (\%). A pixel
  removed for more than one reason is counted once, in the first column that
  applies, from left to right.}
  \label{tab:mask}
  \begin{tabular}{lccccc}
    \toprule
    City & Water & S2 stack & S2 median & AlphaEarth & Total \\
    \midrule
    Milan & 0.78 & 1.82 & 0.00 & 0.00 & 2.61 \\
    Hanoi & 6.52 & 1.51 & 0.00 & 0.00 & 8.03 \\
    HCMC  & 8.94 & 0.00 & 0.00 & 0.00 & 8.94 \\
    \bottomrule
  \end{tabular}
\end{table}

\paragraph{Stratified sampling.} The training target is divided into seven
imperviousness classes: C0 (0\%), C1 (1--20\%), C2 (21--40\%), C3 (41--60\%),
C4 (61--80\%), C5 (81--99\%) and C6 (100\%). Within the valid mask, 500
points are drawn at random from each class, giving a balanced set of
3{,}500 points per city and target (\cref{fig:points}). Stratification
prevents the sample from being dominated by the most frequent values,
typically 0\% and 100\%, and guarantees that intermediate imperviousness is
represented. In Milan, sampling is done separately for CLMS and for GHSL,
since the two products differ pixel by pixel; in Hanoi and HCMC the points
are drawn from GHSL, the only target available.

\begin{figure}[htbp]
  \centering
  \figfrom{final report Fig.~7 --- Milan, the 3{,}500 sample points of each
  training target, coloured by IMD class}
  \caption{Milan: the 3{,}500 sample points of each training target (CLMS,
  GHSL), coloured by imperviousness class.}
  \label{fig:points}
\end{figure}

\paragraph{Predictor extraction.} At every point the predictor values are
extracted once and stored: the 64 AlphaEarth bands, and the Sentinel-2
composite bands (10 for the median, 50 for the percentile composite, and 10
per date for the stack: 40 in Milan, 30 in Hanoi and 10 in HCMC). Each point
thus carries its predictors and its imperviousness value from the training
target, and every model of a city reads the same table.

% ---------------------------------------------------------------------------
\subsection{Spatial train/test split}
\label{sec:split}

No train/test split is made: all 3{,}500 points of a city and training
target are used to tune and to train the model. A test set held out from the
same points would measure agreement with the training product, CLMS or GHSL,
and not with the actual imperviousness on the ground. Accuracy is instead
measured against the independent EarthLabel plots (\cref{sec:validation}),
which are photo-interpreted, never used in training and shared by every map
of a city. The spatial structure of the data is handled within the training
set, through spatially blocked cross-validation (\cref{sec:cv}), which keeps
points from being tested by a model trained on their own neighbourhood.

% ---------------------------------------------------------------------------
\subsection{Spatial-block cross-validation and block-size sensitivity}
\label{sec:cv}

Neighbouring pixels are spatially autocorrelated, so a random
cross-validation would test points against models trained on their
immediate neighbours and overstate skill. The hyperparameters are therefore
chosen by spatially blocked five-fold cross-validation
(\cref{fig:fivefold}). The points are grouped into square spatial blocks,
the blocks are shuffled with a fixed seed and assigned to five folds, and in
each of five rounds the model is trained on four folds and tested on the
fifth.

\begin{figure}[htbp]
  \centering
  \figfrom{final report Fig.~8 --- spatially blocked five-fold
  cross-validation at the three block sizes}
  \caption{Spatially blocked five-fold cross-validation, on a synthetic point
  field, at block sizes of 500, 1{,}000 and 2{,}000\,m. Each block belongs to
  one fold; each round tests on one fold and trains on the other four.}
  \label{fig:fivefold}
\end{figure}

\paragraph{Block-size sensitivity.} Because the appropriate spatial scale of
the blocks is not known in advance, the cross-validation is repeated with
blocks of 500, 1{,}000 and 2{,}000\,m. Block size is treated as part of the
search: every hyperparameter setting is scored at every block size, and the
spread of the scores across block sizes measures how sensitive model
selection is to this choice (reported in \cref{tab:cv}).

\paragraph{Hyperparameter search.} Forty settings are drawn at random
(seed 42) from the grid in \cref{tab:grid}. Each setting is scored at each of
the three block sizes, giving 120 combinations per model, and the
combination with the lowest mean RMSE over the five rounds is selected.
Training and scoring run on the Earth Engine servers with the same random
forest implementation that produces the maps.

\begin{table}[htbp]
  \centering
  \small
  \caption{Random forest hyperparameter grid. Forty settings are drawn at
  random from it; each is scored at three block sizes. $p$ is the number of
  predictor bands.}
  \label{tab:grid}
  \begin{tabular}{ll}
    \toprule
    Hyperparameter & Values \\
    \midrule
    Number of trees & 100, 200, 300, 500 \\
    Bands tried at each split & $\sqrt{p}$, $\log_2 p$, $0.3\,p$, $0.5\,p$ \\
    Minimum points in a leaf & 1, 2, 4 \\
    Maximum nodes per tree & unlimited, 250, 500, 1{,}000 \\
    Block size (cross-validation) & 500, 1{,}000, 2{,}000\,m \\
    \bottomrule
  \end{tabular}
\end{table}

\paragraph{Out-of-fold error.} With the selected setting and block size,
every point is predicted by the model trained without its fold. The RMSE of
these out-of-fold predictions is the cross-validation error of the model.
It is measured against the model's own training target and serves only for
model selection; it is not an accuracy result.

% ---------------------------------------------------------------------------
\subsection{Sentinel-2 compositing strategies}
\label{sec:compositing}

Sentinel-2 Level-2A surface reflectance \cite{esa2018s2,drusch2012sentinel2}
(\texttt{COPERNICUS/S2\_SR\_HARMONIZED}) is used with ten bands: B2, B3, B4,
B5, B6, B7, B8, B8A, B11 and B12. Dates are screened per city by their cloud
percentage over the area of interest and their valid percentage, the share
of pixels left once cloud and no-data pixels are excluded
(\cref{tab:dates}).

\begin{itemize}
  \item \emph{Milan, strict screening} (cloud below 20\%, valid above 90\%):
  30 dates, from 23 March to 8 December 2018. No date in January, February or
  May passed.
  \item \emph{Hanoi and HCMC, loose screening} (cloud below 35\%, valid above
  80\%): four dates for Hanoi and three for HCMC. The looser thresholds are
  needed because few dates in the monsoon climate pass the Milan criteria.
\end{itemize}

\begin{table}[htbp]
  \centering
  \small
  \setlength{\tabcolsep}{5pt}
  \caption{Sentinel-2 dates used per city in 2018, with cloud and valid
  percentage over the area of interest. Dates marked $\ast$ form the stack
  composite of Hanoi and HCMC (the clearest dates of each city).
  % TODO: mark the four Milan stack dates as in the final report's Table 1.
  }
  \label{tab:dates}
  \begin{tabular}{lcc@{\hspace{14pt}}lcc@{\hspace{14pt}}lcc}
    \toprule
    \multicolumn{6}{c}{Milan} & \multicolumn{3}{c}{Hanoi and HCMC} \\
    \cmidrule(lr){1-6}\cmidrule(lr){7-9}
    Date & Cloud & Valid & Date & Cloud & Valid & Date & Cloud & Valid \\
    \midrule
    2018-03-23 & 0.0 & 100.0 & 2018-09-04 & 3.4 & 96.5 & \multicolumn{3}{l}{\emph{Hanoi}} \\
    2018-04-17 & 2.3 & 97.7 & 2018-09-09 & 1.9 & 98.0 & 2018-01-19$^\ast$ & 0.0 & 100.0 \\
    2018-04-22 & 0.0 & 100.0 & 2018-09-14 & 1.8 & 98.2 & 2018-03-10 & 16.8 & 83.2 \\
    2018-06-01 & 3.2 & 96.6 & 2018-09-19 & 6.0 & 94.0 & 2018-10-31$^\ast$ & 0.0 & 100.0 \\
    2018-06-16 & 1.5 & 98.4 & 2018-09-24 & 7.3 & 92.6 & 2018-12-20$^\ast$ & 1.4 & 98.3 \\
    2018-06-21 & 7.8 & 92.1 & 2018-09-29 & 0.3 & 99.7 & \multicolumn{3}{l}{\emph{HCMC}} \\
    2018-07-01 & 4.2 & 95.7 & 2018-10-04 & 0.3 & 99.7 & 2018-02-12$^\ast$ & 0.0 & 100.0 \\
    2018-07-06 & 7.8 & 92.1 & 2018-10-09 & 5.2 & 93.7 & 2018-12-24 & 16.7 & 81.2 \\
    2018-07-26 & 6.0 & 93.9 & 2018-10-14 & 0.8 & 99.2 & 2018-12-29 & 14.2 & 84.7 \\
    2018-07-31 & 1.3 & 98.6 & 2018-10-19 & 2.1 & 97.6 & & & \\
    2018-08-05 & 0.2 & 99.7 & 2018-10-24 & 0.0 & 100.0 & & & \\
    2018-08-10 & 9.1 & 90.7 & 2018-11-18 & 1.6 & 98.2 & & & \\
    2018-08-15 & 1.6 & 98.2 & 2018-11-28 & 1.2 & 98.3 & & & \\
    2018-08-20 & 4.3 & 95.6 & 2018-12-03 & 3.0 & 96.5 & & & \\
    2018-08-30 & 1.5 & 98.5 & 2018-12-08 & 0.0 & 100.0 & & & \\
    \bottomrule
  \end{tabular}
\end{table}

From the screened dates three composites are built for Milan, and two, the
median and the stack, for Hanoi and HCMC. For a given pixel and band, the 30
Milan dates give 30 reflectance values (\cref{fig:onepixel}); the composites
differ in how many of these values they keep.

\begin{itemize}
  \item \emph{Median} keeps one value per band, the middle of the screened
  dates: 10 bands. It is simple and robust to residual cloud, but it hides
  how the pixel changes through the year.
  \item \emph{Percentile} keeps five values per band that describe the
  spread of the dates: the 10th, 25th, 50th, 75th and 90th percentiles,
  50 bands. A pixel that is stable through the year gets five close values; a
  pixel that changes gets widely spread ones. It is built for Milan only,
  where enough dates exist to estimate the percentiles.
  \item \emph{Stack} keeps whole dates unchanged, with no combining: 10
  bands per date. In Milan it uses four near cloud-free dates spread from
  March to October (40 bands), so the model sees each pixel at four moments
  of the year; in Hanoi it uses the three clearest dates (30 bands), and in
  HCMC a single clear date (10 bands). Nothing is gap-filled: a pixel that is
  cloud or shadow on any stack date has no stack value.
\end{itemize}

\begin{figure}[htbp]
  \centering
  \figfrom{final report Fig.~5 --- one pixel, one band: its 30 Milan values,
  median and percentiles}
  \caption{One pixel and one band: its 30 values from the 30 Milan dates
  (dots). The median composite keeps the middle value (p50); the percentile
  composite keeps p10, p25, p50, p75 and p90.}
  \label{fig:onepixel}
\end{figure}

The AlphaEarth predictor needs no compositing: the annual 2018 embedding
(\texttt{GOOGLE/SATELLITE\_EMBEDDING/V1/ANNUAL}) already summarises a full
year of Sentinel-2, Sentinel-1, Landsat and other observations in 64 bands
per 10\,m pixel \cite{brown2025alphaearth}, and is used directly.

% ---------------------------------------------------------------------------
\subsection{Random Forest modelling}
\label{sec:rf}

All models are random forest regressors \cite{breiman2001rf}, implemented
with Earth Engine's \texttt{smileRandomForest} in regression mode. Random
forests handle many correlated predictors, such as the 64 embedding axes or
the 50 percentile bands, without feature selection, and are widely used for
imperviousness and land-cover regression from satellite data.
\litneeded{random forest applications to impervious surface or
imperviousness density estimation}

Each model is tuned as described in \cref{sec:cv}. The final model is then
trained on all 3{,}500 points with the selected setting (bootstrap fraction
1.0, fixed seed). Each model is stored as a record of its training points,
predictor bands, settings and seed, from which Earth Engine rebuilds the
identical forest; the zero-shot maps (\cref{sec:zeroshot}) are produced from
these records.

\paragraph{Raster prediction.} The selected forest is applied to every 10\,m
pixel of the area of interest to produce a continuous imperviousness map,
with water masked. Milan maps are produced in WGS~84 / UTM zone 32N
(EPSG:32632) and the Vietnamese maps in WGS~84 / UTM zone 48N (EPSG:32648).
In Milan there are eight maps: four predictor sets (S2 stack, S2 percentile,
S2 median and AlphaEarth), each trained once on CLMS and once on GHSL
(\cref{tab:maps}). Because the stack is not gap-filled, a stack map has no
value where any of its dates is cloudy.

\begin{table}[htbp]
  \centering
  \small
  \caption{Predicted maps per city. ``Local'' models are trained on points
  from the city itself; ``zero-shot'' maps apply the Milan CLMS-trained model
  unchanged.}
  \label{tab:maps}
  \begin{tabular}{lllc}
    \toprule
    City & Training & Predictors & Maps \\
    \midrule
    Milan & Local, CLMS & S2 stack, S2 percentile, S2 median, AlphaEarth & 4 \\
    Milan & Local, GHSL & S2 stack, S2 percentile, S2 median, AlphaEarth & 4 \\
    Hanoi, HCMC & Local, GHSL & S2 stack, S2 median, AlphaEarth & 3 each \\
    Hanoi, HCMC & Zero-shot (Milan CLMS) & S2 median, AlphaEarth & 2 each \\
    \bottomrule
  \end{tabular}
\end{table}

% ---------------------------------------------------------------------------
\subsection{Local retraining}
\label{sec:local}

In Hanoi and HCMC, a new model is trained for each predictor set on points
from the city itself, following \cref{sec:sampling,sec:cv,sec:rf}: 3{,}500
stratified points drawn from GHSL inside the city's valid mask, predictors
extracted at those points, and a random forest tuned by blocked
cross-validation and trained on all points. Three locally retrained maps are
produced per city, with AlphaEarth, the S2 median and the S2 stack. Local
retraining requires no local reference data beyond the global GHSL product,
so it can be applied to any city; its cost is that the model learns GHSL's
roof-only definition of imperviousness.

% ---------------------------------------------------------------------------
\subsection{Zero-shot cross-city transfer}
\label{sec:zeroshot}

Zero-shot transfer tests whether a model trained in one city can map another
without any local training. The Milan models trained on CLMS, the better of
the two targets, are applied unchanged to Hanoi and HCMC, with no
retraining, recalibration or local points. This is done for two predictor
sets:

\begin{itemize}
  \item \emph{AlphaEarth}: the 64 embedding bands have the same meaning in
  every city, so the Milan model is applied to the Vietnamese 2018
  embeddings directly.
  \item \emph{S2 median}: the Milan model, trained on the median of 30 Milan
  dates, is applied to the median of the three or four screened dates of each
  Vietnamese city, with the same ten bands.
\end{itemize}

The stack has no zero-shot map: its bands are tied to the dates chosen for
each city, and the number of dates differs (four in Milan, three in Hanoi,
one in HCMC), so the Milan stack model cannot be applied elsewhere. The
percentile composite is not built for the Vietnamese cities, which have too
few dates. Together with local retraining, this gives five maps per
Vietnamese city (\cref{tab:maps}).

% ---------------------------------------------------------------------------
\subsection{Independent EarthLabel validation}
\label{sec:validation}

Every map is validated against an independent reference: 450 EarthLabel
plots per city, 1{,}350 in total, photo-interpreted with the EarthLabel
annotation tool (\url{https://github.com/tanxiaoo/earth-label}). No model
sees them during training or tuning.

\paragraph{Plot design.} Each plot is a 10\,m\,$\times$\,10\,m square that
coincides exactly with one pixel of the predicted maps, so each plot is
compared with one map value. The plot is subdivided into nine sub-cells, and
each sub-cell is photo-interpreted as impervious or pervious
(\cref{fig:earthlabel}). The plot value is the share of impervious
sub-cells; two impervious sub-cells out of nine, for example, give 22.2\%.
Plot values therefore take one of ten levels, 0/9 to 9/9.

\begin{figure}[htbp]
  \centering
  \figfrom{final report Fig.~9 --- one EarthLabel plot and its nine
  sub-cells}
  \caption{One EarthLabel plot. The plot covers one 10\,m pixel; its nine
  sub-cells are photo-interpreted one by one, and the share of impervious
  sub-cells becomes the plot value.}
  \label{fig:earthlabel}
\end{figure}

\paragraph{Plot sampling.} Like the training points, the plots are drawn
only inside the valid mask of \cref{sec:sampling}, stratified by the
imperviousness class of the reference map. The plots and the training points
therefore cover the same area, every map has a value at every plot, and all
maps of a city are scored on the same 450 plots.

\paragraph{Labelling.} The ground truth is the finalised interpretation of
two independent labellers, except for 22 replacement plots (6 in Milan, 16
in Hanoi) that were labelled jointly.

\paragraph{Benchmarks.} Besides the predicted maps, the reference products
themselves are scored against the same plots as maps in their own right:
CLMS and GHSL in Milan, and GHSL in Hanoi and HCMC. This places each
predicted map relative to the product it was trained on.

% ---------------------------------------------------------------------------
\subsection{Continuous, binary and ordinal metrics}
\label{sec:metrics}

Each map is compared with the plots in three ways (\cref{fig:techniques}),
which together capture the magnitude, the impervious/pervious separation
and the level structure of the errors.

\begin{figure}[htbp]
  \centering
  \figfrom{final report Fig.~10 --- the three validation techniques A, B
  and C}
  \caption{The three validation techniques: A compares the continuous
  values, B thresholds them at 50\%, C places them on ten ordinal levels.}
  \label{fig:techniques}
\end{figure}

\begin{itemize}
  \item \emph{A, continuous error.} The raw predicted percentage is compared
  with the photo-interpreted percentage, with no threshold: RMSE, MAE, bias
  and \Rsq.
  \item \emph{B, binary at 50\%.} Both values are classified as impervious
  (at or above 50\%) or pervious: overall accuracy (OA) and Cohen's $\kappa$.
  \item \emph{C, ordinal, 10 levels.} Both values are rounded to the nearest
  of the ten plot levels (0/9 to 9/9), with no impervious/pervious threshold:
  exact overall accuracy, overall accuracy within one level (OA $\pm$1), and
  the quadratic weighted kappa (QWK), which counts a two-level error as worse
  than a one-level error.
\end{itemize}

The metrics are defined as
\begin{gather*}
  \mathrm{RMSE} = \sqrt{\frac{1}{n}\sum_{i=1}^{n}(y_i-\hat y_i)^2}, \qquad
  \mathrm{MAE} = \frac{1}{n}\sum_{i=1}^{n}\lvert y_i-\hat y_i\rvert, \\
  \mathrm{Bias} = \frac{1}{n}\sum_{i=1}^{n}(y_i-\hat y_i), \qquad
  R^2 = 1-\frac{\sum_i (y_i-\hat y_i)^2}{\sum_i (y_i-\bar y)^2}, \\
  \kappa = \frac{p_o-p_e}{1-p_e}, \qquad
  \mathrm{QWK} = 1-\frac{\sum_{i,j} w_{ij}O_{ij}}{\sum_{i,j} w_{ij}E_{ij}},
  \qquad w_{ij} = \left(\frac{i-j}{k}\right)^2,
\end{gather*}
where $n$ is the number of plots (450 per city), $y_i$ the photo-interpreted
value, $\hat y_i$ the predicted value and $\bar y$ the mean photo-interpreted
value; $p_o$ and $p_e$ are the observed and chance-expected agreement;
$O_{ij}$ and $E_{ij}$ are the observed and expected numbers of plots at level
$i$ in the reference and level $j$ in the map, and $k=9$. OA is the share of
plots on which map and reference agree; OA $\pm$1 also counts plots whose
levels differ by one.

RMSE, MAE and bias are in percentage points (pp) of imperviousness on the
0--100 scale; lower is better for RMSE and MAE, and zero is best for bias.
Bias is observed $-$ predicted, so a positive bias means the map
under-predicts. \Rsq{} is computed against the 1:1 line, so it penalises a
systematic offset as well as scatter. \Rsq, $\kappa$ and QWK are
dimensionless, and higher is better. To rank the maps of a city, each map is
ranked on \Rsq, $\kappa$ and QWK, and the maps are ordered by their mean
rank.

In addition to the plot-based validation, each predicted map is compared
with its benchmark product over the whole raster, as the mean, mean absolute
and standard deviation of the difference, predicted $-$ benchmark.
